% sections/discussion.tex
% No internal phase labels, model filenames, or run IDs.
% Cross-references use section labels only.

\subsection{Reconciling Three Decades of Contradictory Evidence}

The most important contribution of this work is not a new experimental
observation but a mechanistic resolution of why two incompatible bodies of
evidence about the same biological system have accumulated in parallel for thirty
years without either side being proved wrong. Metcalf and the instructive
tradition~\cite{metcalf1988} were correct at their scale of measurement: in bulk
colony assays with many cells pooled, variance cancels out, mean-field dynamics
prevail, and a sharp cytokine-dependent threshold is the natural result. Hoppe
\etal~\cite{hoppe2016} and Olsson \etal~\cite{olsson2016} were equally correct at
their scale: at single-cell resolution, the same circuit produces fate
heterogeneity that cannot be explained by cytokine dose. Both observations emerge
from the same stochastic model by simply changing the initial molecular
count---a parameter that experimentalists implicitly vary between their assays.
The model therefore resolves the debate by showing that its framing---instructive
\emph{or} stochastic---was incorrect: the system is both, depending on the scale
at which it is observed.

The GCSF/EPO ratio sweep (Section~\ref{sec:ratio_sweep}) adds a second resolution
layer. EPO-only dose-response experiments that find flat commitment probability
are characterising the system in a window where the GCSF/EPO ratio barely
changes---a trivial special case of the full two-cytokine control surface. The
bifurcation revealed by the ratio sweep is not flat: commitment fraction switches
abruptly from greater than 0.92 to zero over a small ratio interval, with no
intermediate values in 500 replicates. This all-or-nothing cliff is the true
bifurcation boundary of the bistable switch, and it depends on the ratio of both
signals, not on either cytokine alone. The clinical implication is direct: EPO
dose escalation in the absence of GCSF normalisation will not cross the
bifurcation threshold; what matters is the relative receptor occupancy of the two
competing signals.

\subsection{The Three-Layer Architecture as a Design Principle}

The decomposition of the commitment cascade into three layers
(Section~\ref{sec:layers}) assigns qualitatively different functional roles to
molecular components whose properties were not previously distinguished
(Table~\ref{tab:layers_disc}). The reception layer sets a coarse erythroid bias
through receptor stoichiometry; it is the system's compass, pointing the cell
toward a fate without specifying it. The decision layer is where the fate is
actually chosen, through molecular-count noise at very small copy numbers; it is
the coin flip embedded in cell biology. The execution layer converts that choice
into a cytokine-maintained chemical record, quickly and EPO-insensitively.

\begin{table}[ht]
\caption{Three-layer architecture as a design principle. Each layer has distinct
EPO-dose sensitivity, noise characteristics, and therapeutic implications.
``Bet-hedging'' refers to population diversity maintained through decision-layer
stochasticity~\cite{losick2008}; ``chromatin record'' refers to the molecular
irreversibility inferred from Pina \etal~\cite{pina2012}.}
\label{tab:layers_disc}
\centering
\small
\begin{tabularx}{\linewidth}{@{}lXXXX@{}}
\toprule
\textbf{Layer} & \textbf{Functional role} & \textbf{EPO-dose sensitivity} &
\textbf{Noise / variability} & \textbf{Therapeutic handle} \\
\midrule
Reception & Sets erythroid bias (identity, not probability) &
EPO-insensitive; set by receptor stoichiometry & Negligible &
Receptor surface residence; EPOR/GCSFR ratio \\
\addlinespace
Decision & Stochastic fate specification & EPO-insensitive &
Dominant (bet-hedging) & GATA1 mRNA stability; basal PU.1 transcription rate \\
\addlinespace
Execution & Cytokine-maintained GATA1 dominance & EPO-insensitive once
triggered & Negligible & GATA1 kinase; PU.1 E3 ubiquitin ligase \\
\bottomrule
\end{tabularx}
\end{table}

This design principle has a broad implication: because all three layers are
EPO-dose-insensitive, the system cannot be finely tuned by adjusting EPO
concentration. EPO functions as a binary permissive signal---it enables
commitment to proceed but does not control its probability or depth. Fine-grained
control over haematopoietic fate must therefore come from elsewhere: from the
copy-number regime of the decision layer (set developmentally), from receptor
stoichiometry (potentially targetable), or from nuclear pH acting on
transcription kinetics that connect all three layers. The cytokine-withdrawal
experiment (Section~\ref{sec:withdrawal}) demonstrates that the execution layer,
without an epigenetic chromatin layer, produces only a cytokine-maintained state,
consistent with the observation of Pina \etal~\cite{pina2012} that an irreversible
molecular record is written before any differentiation marker appears, and with
the bet-hedging principle of Losick and Desplan~\cite{losick2008}.

\subsection{Nuclear pH as a Modulator of Erythroid Commitment Probability}

The EPO--pH sweep (Section~\ref{sec:ph_sweep}) reveals that nuclear pH shifts
the EPO commitment threshold by an amount and in a direction that mean-field
analysis cannot capture. The net-flux prediction of EPO*\,=\,0.52~mM is
pH-independent and underestimates the stochastic threshold at all three pH
values; the pH-dependent shift in EPO* (0.6125 $\to$ 0.5535~mM from pH\,=\,7.0 to
pH\,=\,8.0) is an emergent property of stochastic competition in the decision
layer, not of the kinetic rates alone. This result illustrates a broader
principle: bifurcation thresholds in bistable stochastic circuits depend on
fluctuation statistics as well as mean-field steady states, and environmental
parameters that alter fluctuation amplitude (such as pH acting on inhibition
constants) can shift thresholds in ways invisible to deterministic analysis.

The physiological interpretation of the two pH regimes (increased commitment
probability at alkaline pH near the bifurcation; deeper attractor at acidic pH
in committed cells) bears on the clinical observation that EPO hyporesponsiveness
is most pronounced in patients with chronic kidney disease who also present with
metabolic acidosis. The present model suggests a specific mechanism: at
near-bifurcation EPO concentrations, acidic nuclear pH increases the EPO
barrier to commitment, an effect that EPO dose escalation cannot overcome because
it acts through mutual-inhibition kinetics rather than receptor occupancy.
Conversely, in cells that do commit under acidic conditions, PU.1 silencing is
more complete. The clinical hyporesponsiveness may therefore reflect a reduced
fraction of progenitors committing rather than a defect in those that do
commit---a distinction accessible to joint measurement of commitment fraction and
erythroid marker intensity in pH-manipulated cultures.

The dramatic consequence of the pH shift near the bifurcation---a single EPO
concentration producing 4\% erythroid at pH\,=\,7.0 and 95\% erythroid at
pH\,=\,8.0---also reframes the interpretation of experiments in which nuclear pH
varies across culture conditions, developmental stage, or disease state. Results
from such experiments may report an apparent EPO effect that is partly or wholly
a pH effect.

\subsection{Five Falsifiable Experimental Predictions}

The value of this model for the experimental community lies in generating
non-trivial predictions that would not have been proposed without it. We state
five predictions in decreasing order of conceptual novelty.

The first prediction is the EPO pulse-duration assay. The pulse-duration sweep
(Section~\ref{sec:pulse}) has already self-tested this prediction within the
model, yielding a non-trivial non-monotone result: only the 120~s and 150~s
pulses produce stable post-withdrawal erythroid fate. Both shorter pulses
(insufficient PU.1 erasure) and longer pulses (GATA1 overshoot decaying past the
tipping point) revert to the myeloid attractor. Ghost-attractor reversal
timescale scales systematically with pulse length, a quantitative signature
arising from the larger GATA1 overshoot produced by longer stimulation. The
experimental prediction is fully bracketed: lower bound of the stable erythroid
window between 90 and 120~s; upper bound between 150 and 300~s; confirmed
erythroid-maintaining interval $[120,\,150]$~s. No existing model or framework
has predicted this pulse-length non-monotonicity.

The second prediction addresses the EPO-insensitivity of the execution layer.
Within cells that have committed to the erythroid fate, the kinetics of GATA1
nuclear accumulation after commitment should be identical across EPO doses. Only
the fraction of cells committing varies with EPO dose; for those that do commit,
the depth and speed of the erythroid programme are EPO-dose-independent. This
prediction can be tested by dual-reporter live imaging measuring commitment
kinetics per cell across EPO doses, sorting for committed cells only, and testing
whether GATA1 nuclear accumulation speed and final level are EPO-dose-independent.

The third prediction is the direct demonstration of the EPO* artefact. A clonal
assay in which 1, 10, 50, and 500 cells per well are stimulated across an EPO
dose range should show the apparent commitment threshold sharpening with cell
number and disappearing at the single-cell limit. The model predicts precisely
the shape of this dependence: at single-cell copy numbers the per-cell commitment
probability is a shallow function of EPO dose, so the standard error of the
proportion falls below 5\% only for large $N$. This experiment would directly
demonstrate that three decades of EPO* research has been characterising a
measurement-scale property rather than a cellular mechanism.

The fourth prediction concerns receptor stoichiometry as a dominant control
variable. A titrated EPOR overexpression at fixed EPO dose should change
P(erythroid) more steeply than an equivalent-fold change in EPO concentration.
Doubling EPOR density shifts the EPOR/GCSFR ratio from 4:1 to 8:1---a
qualitatively different signal hierarchy---while doubling EPO concentration shifts
EPOR occupancy by only a few percent. This prediction suggests a rational design
principle for erythropoiesis-stimulating agent development: targeting EPOR
surface residence time rather than ligand concentration.

The fifth prediction is that phosphorylated GATA1 fraction measured at 30 minutes
post-stimulation predicts individual cell commitment outcome better than total
GATA1 protein level. Because phosphorylated GATA1 is the active form in
execution-layer reactions and total GATA1 includes both pools, the phosphorylated
fraction is more mechanistically proximal to the commitment crossover point. A single-cell
imaging mass cytometry experiment measuring both species at 30 minutes and then
tracking final fate at 48 hours would provide a direct test.

\subsection{Drug Discovery Implications}

The three-layer model has direct implications for therapeutic strategies targeting
erythropoiesis. The finding that the execution layer is EPO-dose-insensitive
mechanistically explains a clinical observation that has lacked a molecular basis:
EPO hyporesponsiveness in myelodysplastic syndrome, chronic kidney disease, and
chemotherapy-induced anaemia, where EPO dose escalation beyond a certain point
does not proportionally increase reticulocyte production. If commitment fraction
is bounded by decision-layer noise set by intracellular molecular copy numbers
rather than EPO concentration, the ceiling cannot be raised by cytokine dose
escalation. Rational therapy for EPO hyporesponsiveness must therefore address
the decision layer---stabilising GATA1 mRNA, reducing GATA1 protein turnover, or
reducing the basal PU.1 transcription rate---rather than the reception layer.

The execution layer (GATA1 kinase and PU.1 E3 ubiquitin ligase) represents a
class of drug targets not previously highlighted in this context. These enzymes
activate within 48 seconds and are EPO-dose-insensitive, functioning as
autonomous commitment amplifiers once engaged. Inhibiting GATA1 phosphorylation
would prevent execution-layer engagement even in cells where the decision layer
had begun to tip toward erythroid---a distinct mechanism from targeting upstream
transcription or receptor signalling. Similarly, PU.1 stabilisation through
deubiquitinase inhibitors would maintain a myeloid-competent state even in the
presence of EPOR signalling, an approach relevant to forced myeloid
differentiation in AML therapy.

\subsection{Limitations}

Several limitations must be stated clearly. The model makes quantitative
predictions based on calibrated parameters, but the parameter values are not
individually measured from a specific cell line; they are derived from published
estimates across multiple experimental systems and adjusted for thermodynamic
consistency. The model should be considered a hypothesis-generating tool rather
than a quantitatively precise replica of any specific cell type. In particular,
the predicted priming-crossover time of 146~s is a model-specific number that
could shift by a factor of two with different parameter choices; the qualitative
prediction---that there exists an early cytokine-dependent dominance point---is
more robust than the specific timing.

The model does not include a chromatin remodelling layer. The cytokine-withdrawal
experiment demonstrates that the erythroid attractor established by EPOR
signalling is reversible: EPO removal produces complete fate reversal within
approximately 1500~s. True transcriptional irreversibility in vivo is mediated by
histone methylation, chromatin accessibility remodelling, and locus-specific DNA
methylation. Adding an epigenetic layer is the highest-priority structural
extension of this model.

The commitment threshold used to classify replicates
($\text{GATA1}_\text{nuc} > 1.5 \times \text{PU.1}_\text{nuc}$) is motivated by the bistable
structure of the model but has not been calibrated against a specific
differentiation marker. Translation from model outcome to a measurable marker
requires additional calibration against experimental data.

The model does not yet reproduce the distribution of commitment times observed in
live-cell imaging experiments. Hoppe \etal~\cite{hoppe2016} measured the time from
GATA1 reporter expression onset to binary fate segregation and observed a
distribution spanning several hours. Quantitatively reproducing this distribution
would constitute the most stringent validation and is the highest-priority
empirical next step.

\subsection{Relationship to Prior Computational Work}

The present model differs from Huang \etal~\cite{huang2007} and Chickarmane
\etal~\cite{chickarmane2009} in three key respects: it includes stochastic
dynamics rather than deterministic ordinary differential equations; it integrates
the full signalling cascade from receptor to nuclear phosphorylation rather than
abstracting the circuit to a two-equation toggle; and it identifies the execution
layer as structurally distinct from the decision layer by its EPO-dose-insensitivity.
The layer decomposition does not appear in prior literature and constitutes the
principal analytical contribution of this work. The localisation of stochastic
noise specifically to the decision layer, and not to the other layers, is a new
claim that is testable experimentally.
